
\chapter*{General Introduction}
\markboth{General Introduction}{}
\addcontentsline{toc}{chapter}{General Introduction}

In a world where nearly every aspect of life depends on interconnected systems, modern society has become heavily reliant on digital technologies such as communication networks, banking systems, healthcare platforms, transportation systems, and even critical infrastructures.

However, this growing dependence has turned connectivity into a vulnerability to be exploited through cyber attacks, making cybersecurity a necessity.

And in the face of ever-growing and increasingly sophisticated threats, a race began.
Among the many security mechanisms developed to address these threats, Network Intrusion Detection Systems (NIDS) play a critical role.

As the name suggests, NIDS operate at the network level, acting as a key line of defense by monitoring network traffic and detecting malicious activities before they can further propagate into the system.

However, the effectiveness of traditional NIDS has become increasingly challenged by the evolving nature of cyber threats. Built mostly around predefined rules and signatures, these systems remain rigid and reactive, making it difficult for them to keep pace with rapidly changing attack patterns.

Another major concern is reaction time, which remains one of the most critical factors in cybersecurity. The impact of an attack is closely tied not only to how quickly it is detected, but also to how quickly it can be analyzed and responded to. Any delay in this process increases the window of exposure and can significantly amplify its overall impact.  Therefore, minimizing response time becomes essential, something traditional NIDS do not always effectively provide

With the rise of artificial intelligence and its integration into almost every domain, AI emerged as a promising solution to many complex problems, including cybersecurity. By introducing learning and adaptability capabilities.

In this context, Our project focuses on the design and implementation of an AI-based Network Intrusion Detection System.

Throughout this project, we adopt the \textbf{CRISP-DM} methodology \cite{CRISP}.

The first chapter, entitled \textbf{<<Project Context>>}, is devoted to the presentation of the host organization, the analysis of existing solutions, and the definition of the problem statement with the proposed solution.

The second chapter, entitled \textbf{<<Requirements Specification>>}, is dedicated to defining the functional and non-functional requirements, identifying the actors, and modeling the use cases.

The third chapter, \textbf{<<Conceptual Design>>}, presents the conceptual design of the system through the main UML diagrams \cite{UML}.

The fourth chapter, \textbf{<<Data Collection and Preparation>>}, covers data acquisition, data understanding, and data preparation according to the Data Understanding and Data Preparation phases of the CRISP-DM methodology.

The fifth chapter, \textbf{<<Modeling and Evaluation>>}, describes the adopted modeling approach, the training process, and the evaluation of model performance.

The sixth chapter, \textbf{<<Realization>>}, presents model integration, deployment architecture, and the main interfaces of the solution.

Finally, we conclude this report with a general conclusion that summarizes the work carried out, the main contributions, and possible future improvements.